\documentclass[10pt,oneside]{article}

%%%%%%%%%%%%%
\setlength{\textheight}{8.75in} %Letter is 11in, less 2 for margins, less 0.25 for footer
\setlength{\oddsidemargin}{0.0in} %gets +1inc
\setlength{\evensidemargin}{0.0in} %gets +1inch
\setlength{\textwidth}{6.50in} %Letter is 8.5, less 2 inches for margins
\setlength{\topmargin}{0.5in}
\setlength{\headheight}{0in}
\setlength{\headsep}{0in}
\setlength{\parindent}{0.25in}
%%%%%%%%%%%%

% use letters instead of symbols to accomodate >7 authors
\makeatletter
\let\@fnsymbol\@alph
\makeatother

\usepackage[utf8]{inputenc}
\usepackage[numbers]{natbib}
\usepackage{graphicx}
\usepackage[colorlinks=true,citecolor=black,urlcolor=blue]{hyperref}

\title{%Hack to get the logo on the PDF front page:
\vspace{-1.5in}
\includegraphics[width=0.4\textwidth]{biopython.jpg}\\
\vspace{3mm}Biopython Project Update 2016}
\author{
	\underline{Christian Brueffer}\thanks{Department of Clinical Sciences, Lund University, Lund, SE. Email: \href{mailto:christian.brueffer@med.lu.se}{christian.brueffer@med.lu.se}},
    Tiago Ant\~{a}o\thanks{Division of Biological Sciences, University of Montana, Montana, USA},
    Peter Cock\thanks{Information and Computational Sciences, James Hutton Institute (formerly SCRI), Invergowrie, Dundee, UK},
    Eric Talevich\thanks{Department of Dermatology, University of California San Francisco, San Francisco, CA, USA},
    Michiel de Hoon\thanks{Division of Genomic Technologies, RIKEN Center for Life Science Technologies, Yokohama, JP},
	\\
    Wibowo Arindrarto\thanks{Sequencing Analysis Support Core, Leiden University Medical Center, Leiden, NL},
    Leighton Pritchard\textsuperscript{c},  % same as Peter, saves space
    Anuj Sharma\thanks{Department of Informatics and Telecommunications, University of Athens, Athens, Greece},
    Eric Rasche\thanks{Center for Phage Technology, Department of Biochemistry and Biophysics, Texas A\&M University, USA},
    Aaron Rosenfeld\thanks{Department of Biomedical Engineering, Science \& Health Systems, Drexel University, USA},
    \\
    Connor T. Skennerton\thanks{Division of Geological and Planetary Sciences, California Institute of Technology, Pasadena, CA, USA},
    Marco Galardini\thanks{EMBL-EBI, Wellcome Trust Genome Campus, Cambridge, UK},
    Markus Piotrowski\thanks{Department of Plant Physiology, Ruhr-Universität Bochum, DE},
    and the Biopython Contributors}
\date{17\textsuperscript{th} Bioinformatics Open Source Conference (BOSC) 2016, Orlando, USA}

\begin{document}
\maketitle
\thispagestyle{empty}

\vspace{-0.2in}
\noindent
Website: \url{http://biopython.org} \\
Repository: \url{https://github.com/biopython/biopython} \\
License: Biopython License Agreement (MIT style, see \url{http://www.biopython.org/DIST/LICENSE}) \\

The Biopython Project is a long-running distributed collaborative effort,
supported by the Open Bioinformatics Foundation, which develops a freely
available Python library for biological computation \cite{AppNote}.

We present here details of the latest Biopython release - version 1.66. New
features include: extended Bio.KEGG and Bio.Graphics modules to support drawing
KEGG pathways with transparency; extended ``abi'' Bio.SeqIO parser to decode
almost all documented fields used by ABIF instruments; a QCPSuperimposer
module using the Quaternion Characteristic Polynomial algorithm for superimposing
structures to Bio.PDB; and an extended Bio.Entrez module to implement the NCBI
Entrez Citation Matching function and to support NCBI XML files with XSD schemas.
Additionally we fixed miscellaneous bugs, enhanced our test suite and continued our
efforts to abide by the PEP8 coding style guidelines.

We are currently preparing a new release -- version 1.67 -- that will deprecate the
ability to compare SeqRecord objects with ``=='', which sometimes lead to surprising
results.  In addition it will feature a new experimental Bio.phenotype module for
working with Phenotype Microarray data; updates to Bio.Data to include NCBI genetic
code table 25, covering Candidate Division SR1 and Gracilibacteria; an update to
Bio.Restriction to include the REBASE May 2016 restriction enzyme list; updates to
BioSQL to use foreign keys with SQLite3 databases; as well as corrections
to the Bio.Entrez module and the MMCIF structure parser.

Our continuous integration process on GitHub has been enhanced by including
external services like Landscape, Quantified Code and Codecov to perform
quality review, test coverage analysis and generation of quality metrics.

Finally, our range of Docker containers has been greatly enhanced. In addition to
a basic container that includes Python 2 and 3 with Biopython and all its
dependencies, as well as a BioSQL container, we now also provide two versions
of Jupyter notebook containers: a basic one, and a version including the
Biopython tutorial as notebooks.


\begin{thebibliography}{}

\bibitem[Cock {\it et al}., 2009]{AppNote}Cock, P.J.A., Antao, T., Chang, J.T., Chapman, B.A., Cox, C.J., Dalke, A., Friedberg, I., Hamelryck, T., Kauff, F., Wilczynski, B., de Hoon, M.J. (2009) Biopython: freely available Python tools for computational molecular biology and bioinformatics. {\it Bioinformatics} {\bf 25}(11) 1422-3. \href{http://dx.doi.org/10.1093/bioinformatics/btp163}{doi:10.1093/bioinformatics/btp163}

\end{thebibliography}

\end{document}
